\begin{solapartat}
L'àrea del circuit tancat que formen la forca i la vareta en funció del temps és:
\[ \text{Àrea} = L\,v\,t \ \text{\punts{0,2}} \]
El flux del camp magnètic que passa per aquesta àrea serà:
\[ \Phi = B\,L\,v\,t \ \text{\punts{0,2}} \]
La força electromotriu generada en el circuït serà:
\[ \epsilon = \left|\frac{\mathrm{d}\Phi}{\mathrm{d}t}\right| = B\,L\,v \ \text{\punts{0,2}} = 3\ \mathrm{V} \ \text{\punts{0,2}} \]
El sentit de la circulació del corrent serà el contrari que tindria si el mateix corrent hagués de crear el camp magnètic, per tant serà en sentit horari \punts{0,2}
\end{solapartat}

\begin{solapartat}
A partir de la llei de Ohm tindrem:
\[ I = \frac{\epsilon}{R} = \frac{3\mathrm{V}}{30\Omega} = 0,1\ \mathrm{A} \ \text{\punts{0,2}} \]
La força que fa el camp magnètic sobre la vareta serà:

\figura[0.45]{sol_fig1.png}
\aladreta{\punts{0,2}}

\[ \vec{F} = L\,\vec{I}\wedge\vec{B} \ \text{\punts{0,2}} \Rightarrow |\vec{F}| = L\,I\,B = 0,05\ \mathrm{N} \ \text{\punts{0,2}} \]
Per tant, per tal que la vareta segueixi amb velocitat constant, haurem de fer una força igual i de sentit contrari a la trobada anteriorment. \punts{0,2}
\end{solapartat}
